% ============================================================
% M3 S3 练习件·W3 臂E —— 课时10 曲线与方程（2.4）＝照抄母版件（课时01）体例
% 生成：M3 成卷轮 S3 W3 臂E｜2026-09-14｜编译 cwd＝本目录（中文路径禁 \input 绝对路径）
%   xelatex main-true.tex ×2 → 含详解印本；xelatex main-false.tex ×2 → 纯题版（单源双本）
% 输入链（全只读）：题面库 课时10-2.4曲线与方程.md（题面逐字装配）＋ -答案侧.md（值逐字回嵌）
%   ＋ 定稿/批C-课时10-2.4曲线与方程.md（详解回嵌源）；toolchain＝qp-m3.sty（钉后版
%   md5 7c3930362be8a0a2bdf21bbf8ac16573，本地挂载副本，破损即停）。
% 母版体例：详见《练习件母版体例说明.md》（课时01 件）；门谱读数见 过程件 工作区/_tmpM3S3W3臂E0914/。
% 键账：件内 ansblock 键集 ≡ 题面库片练键集 ≡ manifest 练键序 简1~简10＋中1~中4＋难1~难2
%   （16 键，零缺漏零多余）；拓区隔离：2章-拓/测/滚 键不入本件（S4 拓展册收；
%   10-拓1 撤位席不设键照 canonical 总表 §四.2），导学键 G1~G5 属导学件域，亦不入。
% 卷式例外案B：练习件不属卷式——答案块物理落点恒为题后 ansblock，禁卷末附卷制；
%   全线括线模（导言 \ansblockgrayfalse 一次置定，件内不切换；灰底盒不可跨栏断，练习件禁用）。
% 题序＝manifest 键序（零重排）；印面题号 1~16 连号（\tihao 号＝\ansitem 号同键）；
%   三档切位 1~10／11~14／15~16（简10／中4／难2 照库序切位，非难度重排）。
% 槽配实录：单选3（简7/简8/难1）＋多选0＋填空1（难2）＋解答12（简1~简6/简9/简10/中1~中4）；
%   10选4填2解 为波次方案基线槽配，本片实配照题面侧头标（批C §三 10简4中2难 同口径；
%   10-难1「多结论选择」源卷未标多选，manifest 多选＝0，不计多选门、零 \duoxuan）。
% 0913 批C置换入槽在卷：简8/简9/简10＝命制 10-命1改/命2/命3改（题面/值/详解照录）；
%   原中5~7 三席挤出移拓展册（10-拓13~15），拓1 撤席让位——域内零影响，件内照 16 练键装配。
% 选项槽位：默认四连排（0.25\linewidth−0.25\lxhang），超宽降档梯 四连排→两连排
%   （0.5\linewidth−0.5\lxhang−0.25em）→单列 \lxopt；禁缩字号/负 kern/删标点凑宽。
%   本片实测降档：见值台账「降档登记」＋回执（槽宽门真算读数）。
% 解答书写区：简1~简6/简9/中2~中4 无小问 16mm；简10 三小问 32mm；中1 两问 24mm
%   （16~40mm 带内；纯题档吞 ansblock 不吞 \liubai）。
% 填空印答：难2 空位 \kongda{②④}（详解版印答、纯题版退定宽空线，两档等形）。
% 难度词：简→简单、0.65（中）→中档（批C 派生同批B 口径）、难→难照录头标；知识点 N 照
%   导学件10 知识导学三分册：一＝曲线的方程与方程的曲线、二＝求轨迹方程·直接法、
%   三＝由方程研究曲线（派生登记见值台账）。
% 印面清洗：题面侧/定稿【注】行不入印面；亲算落账行不入印面；详解 ✓ 记号剔（检验句文句
%   保留）；题面 ASCII 引号照导学件10 同题先例宏化为 CJK 弯引号；ASCII 箭头宏化 \(\rightarrow\)。
% 图注记：难1 题面〔图注〕照录自源（图件由成卷轮图形工位统一出，图债清单见回执，
%   占位零漂不阻塞编译门）。
% ============================================================
\documentclass[fontset=none]{ctexart}
\usepackage{qp-m3}
% —— 印面逐字符号路由补钉（承导学件母版；− 走 TNR、▱ NSC 子块；禁改 sty 保 md5） ——
\xeCJKDeclareCharClass{Default}{"2212}%
\setCJKfamilyfont{nscblk}[Path=C:/提示词/工作区/字替对照-0909/variantF/fonts/,BoldFont=NSC-w600.ttf]{NSC-w500.ttf}%
\xeCJKDeclareSubCJKBlock{nscblk}{"25B1}%
\newcommand{\pxparallelogram}{{\CJKfamily{nscblk}▱}}%
% —— 断行微调（承导学件母版实测档，三零门承重；\emergencystretch 不另设） ——
\xeCJKsetup{CJKglue={\hskip 0.10em plus 0.08em minus 0.05em}}%
% —— 练习件全线括线模（S3 波次方案 §四.3：导言一次置定、件内不切换） ——
\ansblockgrayfalse
% —— 页脚件名（qp-headfoot 复用入口） ——
\renewcommand{\qpzhangming}{第二章\quad 平面解析几何}%
\renewcommand{\qpjianming}{练习件}%

\begin{document}
%装配轮S6三本跨本连续页码（G7方案B）setcounter 注入位
\setcounter{page}{199}

\keshi{第10课时\quad 曲线与方程}
\begin{multicols}{2}
\emergencystretch=1em

% ============================================================
% 基础巩固（题1~10＝简1~简10：解答9＋单选2；题后紧跟答案制，全线括线 ansblock）
% ============================================================
\dangbadge{基}{础}{巩}{固}

\tihao{1}\zu{点与曲线对应判断×解答} \tieside{简单(知识点一)}判断\(P(\sqrt{3},0)\)，\(Q(-2,3)\)是否在方程\(4x^{2}+3y^{2}=12\)的曲线上．
\liubai[16mm]{解答书写区}

\begin{ansblock}[2章-练-课时10-简1]
% ans:2章-练-课时10-简1
\ansitem{1}{P在曲线上；Q不在曲线上}
\ansnote{详解}{曲线上的点的坐标必为方程的解，反之方程的解对应曲线上的点，故代入即判。\(P(\sqrt{3},0)\)：\(4\times(\sqrt{3})^{2}+3\times 0^{2}=12+0=12\)，是方程的解，\(P\)在曲线上。\(Q(-2,3)\)：\(4\times(-2)^{2}+3\times 3^{2}=16+27=43\neq 12\)，不是方程的解，\(Q\)不在曲线上。}
\end{ansblock}

\tihao{2}\zu{曲线过点求参数×解答} \tieside{简单(知识点一)}已知方程\((x+2)^{2}+(y-3)^{2}=r^{2}\)的曲线通过点\((-1,5)\)，求\(r^{2}\)的值．
\liubai[16mm]{解答书写区}

\begin{ansblock}[2章-练-课时10-简2]
% ans:2章-练-课时10-简2
\ansitem{2}{r²=5}
\ansnote{详解}{曲线通过点\((-1,5)\)，即\((-1,5)\)的坐标是方程的解：\((-1+2)^{2}+(5-3)^{2}=r^{2}\)，得\(r^{2}=1+4=5\)。}
\end{ansblock}

\tihao{3}\zu{纯粹性/完备性说理×解答} \tieside{简单(知识点一)}到两坐标轴距离相等的点组成的轨迹的方程是\(x-y=0\)吗？为什么？
\liubai[16mm]{解答书写区}

\begin{ansblock}[2章-练-课时10-简3]
% ans:2章-练-课时10-简3
\ansitem{3}{不是。x−y=0 只表示其中一条直线，轨迹应为 x−y=0 或 x+y=0（即|x|=|y|）；以 x−y=0 为轨迹方程不完备（漏了直线 x+y=0 上的点）}
\ansnote{详解}{设动点\(M(x,y)\)。到\(x\)轴距离为\(|y|\)、到\(y\)轴距离为\(|x|\)，条件\(|x|=|y|\Longleftrightarrow x-y=0\)或\(x+y=0\)，轨迹是两条互相垂直的直线。方程\(x-y=0\)的解对应的点都在轨迹上（纯粹性成立），但轨迹上的点（如\((1,-1)\)）不都是该方程的解（完备性不成立）——它只描述第一、三象限角平分线，漏了第二、四象限角平分线。故“轨迹的方程是\(x-y=0\)”不对。}
\end{ansblock}

\tihao{4}\zu{直接法·垂直平分线轨迹×解答} \tieside{简单(知识点二)}设\(A(-1,-1)\)，\(B(3,7)\)，求线段\(AB\)的垂直平分线的方程．
\liubai[16mm]{解答书写区}

\begin{ansblock}[2章-练-课时10-简4]
% ans:2章-练-课时10-简4
\ansitem{4}{x+2y−7=0}
\ansnote{详解}{设\(P(x,y)\)为线段\(AB\)垂直平分线上任意一点，则\(|PA|=|PB|\)：\((x+1)^{2}+(y+1)^{2}=(x-3)^{2}+(y-7)^{2}\)，展开：\(x^{2}+2x+1+y^{2}+2y+1=x^{2}-6x+9+y^{2}-14y+49\)，化简得\(8x+16y-56=0\)，即\(x+2y-7=0\)。反向检验：满足\(x+2y-7=0\)的点均有\(|PA|=|PB|\)，在\(AB\)的垂直平分线上（两性完备）。验算：\(AB\)中点\((1,3)\)：\(1+6-7=0\)；\(k_{AB}=(7+1)/(3+1)=2\)，所求直线斜率\(-1/2\)，\(2\times(-1/2)=-1\)。}
\end{ansblock}

\tihao{5}\zu{定比轨迹×解答} \tieside{简单(知识点二)}求到两定点\(A(-1,2)\)，\(B(3,2)\)的距离之比为\(\sqrt{2}\)的点的轨迹方程．
\liubai[16mm]{解答书写区}

\begin{ansblock}[2章-练-课时10-简5]
% ans:2章-练-课时10-简5
\ansitem{5}{(x−7)²+(y−2)²=32}
\ansnote{详解}{设\(M(x,y)\)，\(|MA|/|MB|=\sqrt{2}\)，即\(|MA|^{2}=2|MB|^{2}\)：\((x+1)^{2}+(y-2)^{2}=2[(x-3)^{2}+(y-2)^{2}]\)。展开：\(x^{2}+2x+1+(y-2)^{2}=2x^{2}-12x+18+2(y-2)^{2}\)，移项整理：\(x^{2}-14x+(y-2)^{2}+17=0\)，配方得\((x-7)^{2}+(y-2)^{2}=32\)。轨迹为以\((7,2)\)为圆心、\(4\sqrt{2}\)为半径的圆（阿波罗尼斯圆雏形，此名可不向学生出现）。验算：圆与直线\(AB\)（\(y=2\)）交于\(M_{1}(7-4\sqrt{2},2)\)、\(M_{2}(7+4\sqrt{2},2)\)。对\(M_{1}\)：\(|M_{1}A|=|7-4\sqrt{2}+1|=8-4\sqrt{2}=4(2-\sqrt{2})\)，\(|M_{1}B|=|4\sqrt{2}-4|=4(\sqrt{2}-1)\)（取绝对值），比值\(=(2-\sqrt{2})/(\sqrt{2}-1)=\sqrt{2}(\sqrt{2}-1)/(\sqrt{2}-1)=\sqrt{2}\)；对\(M_{2}\)同法比值亦\(\sqrt{2}\)。}
\end{ansblock}

\tihao{6}\zu{轨迹方程·斜率积×解答} \tieside{简单(知识点二)}已知\(\triangle ABC\)底边两端点\(B(0,6)\)、\(C(0,-6)\)，若这个三角形另外两边所在直线的斜率之积为\(-4/9\)，求点\(A\)的轨迹方程．
\liubai[16mm]{解答书写区}

\begin{ansblock}[2章-练-课时10-简6]
% ans:2章-练-课时10-简6
\ansitem{6}{x²/81+y²/36=1（x≠0）}
\ansnote{详解}{设\(A(x,y)\)，\(x\neq 0\)（\(x=0\)时\(A\)在\(y\)轴上与\(B\)、\(C\)共线，不能成三角形且斜率不存在）。\(k_{AB}\cdot k_{AC}=[(y-6)/x]\cdot[(y+6)/x]=(y^{2}-36)/x^{2}=-4/9\)，去分母：\(9(y^{2}-36)=-4x^{2}\)，整理得\(4x^{2}+9y^{2}=324\)，即\(x^{2}/81+y^{2}/36=1\)（\(x\neq 0\)）。}
\end{ansblock}

\tihao{7}\zu{方程与曲线对应辨析×单选} \tieside{简单(知识点一)}下列各组方程中，表示同一条曲线的是\nobreak(\kongwei)

\lxopt{A．\(y=|x|\) 与 \(y=\sqrt{x^{2}}\)}

\lxopt{B．\(y=x\) 与 \(y=\sqrt{x^{2}}\)}

\lxopt{C．\(|y|=x\) 与 \(y=\sqrt{x^{2}}\)}

\lxopt{D．\(x^{2}+y^{2}=0\) 与 \(y=0\)}

\begin{ansblock}[2章-练-课时10-简7]
% ans:2章-练-课时10-简7
\ansitem{7}{A}
\ansnote{详解}{A：\(\sqrt{x^{2}}=|x|\)，两方程解集（图象）完全相同，表同一条曲线；B：\(y=x\)为直线，\(y=\sqrt{x^{2}}=|x|\)为折线（\(x<0\)时\(y=-x\)），异；C：\(|y|=x\Longleftrightarrow x\geqslant 0\)且\(y=\pm x\)，为两条射线，而\(y=|x|\)还含\(x<0\)部分，异；D：\(x^{2}+y^{2}=0\Longleftrightarrow\)点\((0,0)\)，\(y=0\)为整条\(x\)轴，异。}
\end{ansblock}

\tihao{8}\zu{判对应辨析·定点式×单选} \tieside{简单(知识点一)}与方程 \(x^{2}+y^{2}=4\) 表示同一条曲线的是\nobreak(\kongwei)
\bindopt {\fontsize{10.5pt}{\qpTJgLeadLiti}\selectfont \makebox[\dimexpr0.5\linewidth-0.5\lxhang-0.25em\relax][l]{A．\(y=\sqrt{4-x^{2}}\)}\makebox[\dimexpr0.5\linewidth-0.5\lxhang-0.25em\relax][l]{B．\((x-2)^{2}+y^{2}=4\)}\par}
\bindopt {\fontsize{10.5pt}{\qpTJgLeadLiti}\selectfont \makebox[\dimexpr0.5\linewidth-0.5\lxhang-0.25em\relax][l]{C．\(x^{2}+y^{2}-4=0\)}\makebox[\dimexpr0.5\linewidth-0.5\lxhang-0.25em\relax][l]{D．\(|x|+|y|=2\)}\par}

\begin{ansblock}[2章-练-课时10-简8]
% ans:2章-练-课时10-简8
\ansitem{8}{C}
\ansnote{详解}{C与原方程仅差移项，同解变形，解集恒等，表同一条曲线；A：\(\sqrt{4-x^{2}}\)要求\(y\geqslant 0\)，仅为上半圆，缺\(y<0\)部分不完备，排除；B：圆心\((2,0)\)、半径2的圆，与圆心在原点的圆非同一条曲线，排除；D：\(|x|+|y|=2\)是以\((2,0)\)、\((0,2)\)、\((-2,0)\)、\((0,-2)\)为顶点的正方形边界，点\((1,1)\)在其上却不在圆上，排除。故选C。}
\end{ansblock}

\tihao{9}\zu{直接法基础轨迹×解答} \tieside{简单(知识点二)}已知动点 \(M\) 到 \(x\) 轴的距离与到定点 \(A(0,4)\) 的距离相等，求 \(M\) 的轨迹方程。
\liubai[16mm]{解答书写区}

\begin{ansblock}[2章-练-课时10-简9]
% ans:2章-练-课时10-简9
\ansitem{9}{x²=8(y−2)}
\ansnote{详解}{设\(M(x,y)\)。由题\(|y|=\sqrt{x^{2}+(y-4)^{2}}\)。两边平方（两侧非负，等价）：\(y^{2}=x^{2}+y^{2}-8y+16\)，即\(x^{2}=8(y-2)\)。检验：方程上任一点均满足定义（平方等价可逆），无杂点；\(A(0,4)\)代入\(16\neq 0\)不在轨迹上，无需剔除。故轨迹方程为\(x^{2}=8(y-2)\)。}
\end{ansblock}

\tihao{10}\zu{方程与曲线概念判断×解答} \tieside{简单(知识点一)}判断下列命题真假并说明理由：
(1) 方程\(x^{2}+y^{2}=0\)表示一条直线；
(2) \(|x|=|y|\)的轨迹方程是\(y^{2}=x^{2}\)；
(3) 若点\(P(x_{0},y_{0})\)的坐标是曲线\(C\)的方程\(f(x,y)=0\)的一组解，则\(P\)在\(C\)上。
\liubai[32mm]{解答书写区}

\begin{ansblock}[2章-练-课时10-简10]
% ans:2章-练-课时10-简10
\ansitem{10}{(1) 假　(2) 真　(3) 真}
\ansnote{详解}{(1)\(x^{2}+y^{2}=0\)当且仅当\(x=y=0\)，表示原点一个点，非直线；(2)\(|x|=|y|\)当且仅当\(x^{2}=y^{2}\)（两侧均非负，平方同解），即\(y^{2}=x^{2}\)，命题真；(3)由曲线方程定义的完备性（以方程的解为坐标的点都在曲线上），命题真。}
\end{ansblock}

% ============================================================
% 综合提升（题11~14＝中1~中4：解答4；新拟定稿题全文照 §4 装配）
% ============================================================
\dangbadge{综}{合}{提}{升}

\tihao{11}\zu{由方程研究曲线×解答} \tieside{中档(知识点三)}已知曲线\(C\)：\(x^{2}+y^{2}-2|x|=0\)。(1)指出\(C\)的形状；(2)写出\(C\)的对称性、范围、与坐标轴的交点，并求\(C\)围成的区域的面积。
\liubai[24mm]{解答书写区}

\begin{ansblock}[2章-练-课时10-中1]
% ans:2章-练-课时10-中1
\ansitem{11}{(1)两个外切于原点的单位圆（「8」字形）；(2)关于x轴、y轴、原点均对称；x∈[−2,2]，y∈[−1,1]；与y轴仅交于(0,0)，与x轴交于(±2,0)与(0,0)；S=2π}
\ansnote{详解}{(1)按\(|x|\)分段：\(x\geqslant 0\)时\((x-1)^{2}+y^{2}=1\)，\(x<0\)时\((x+1)^{2}+y^{2}=1\)，两圆圆心\((\pm 1,0)\)、半径1，外切于原点。\par\noindent (2)\(|x|\)使\(C\)关于\(y\)轴对称；两圆连心线在\(x\)轴上\(\rightarrow\)关于\(x\)轴对称；由两者\(\rightarrow\)关于原点对称。范围\(x\in[-2,2]\)、\(y\in[-1,1]\)。\(x=0\Longrightarrow y^{2}=0\)仅\((0,0)\)；\(y=0\Longrightarrow|x|=2\)或\(0\)，得\((\pm 2,0)\)与\((0,0)\)。两圆仅切于一点无重叠，\(S=2\cdot\pi\cdot 1^{2}=2\pi\)。}
\end{ansblock}

\tihao{12}\zu{轨迹方程＋纯粹性完备性检验×解答} \tieside{中档(知识点二)}已知\(A(-2,0)\)、\(B(2,0)\)，动点\(P\)满足\(k_{PA}\cdot k_{PB}=-1/4\)，求\(P\)的轨迹方程，并验证轨迹的纯粹性与完备性。
\liubai[16mm]{解答书写区}

\begin{ansblock}[2章-练-课时10-中2]
% ans:2章-练-课时10-中2
\ansitem{12}{x²/4+y²=1（x≠±2）}
\ansnote{详解}{设\(P(x,y)\)。斜率存在要求\(x\neq\pm 2\)；\(y^{2}/(x^{2}-4)=-1/4\Longrightarrow 4y^{2}=4-x^{2}\)，即\(x^{2}/4+y^{2}=1\)（\(x\neq\pm 2\)）。纯粹性：轨迹上任一点的坐标推导步步可逆，都满足斜率方程；完备性：满足题设的点都满足\(x^{2}/4+y^{2}=1\)且\(x\neq\pm 2\)，都在轨迹上（\(y=0\)时斜率积为\(0\neq -1/4\)，端点自然排除）。故轨迹为椭圆\(x^{2}/4+y^{2}=1\)除去长轴两端点\((\pm 2,0)\)。}
\end{ansblock}

\tihao{13}\zu{轨迹存在性×解答} \tieside{中档(知识点三)}已知圆\(C_{1}\)：\(x^{2}+(y-m)^{2}=9\)与圆\(C_{2}\)：\(x^{2}+(y+2)^{2}=1\)有公共点，求实数\(m\)的取值范围。
\liubai[16mm]{解答书写区}

\begin{ansblock}[2章-练-课时10-中3]
% ans:2章-练-课时10-中3
\ansitem{13}{m∈[−6,−4]∪[0,2]}
\ansnote{详解}{圆心\(C_{1}(0,m)\)、\(C_{2}(0,-2)\)，半径3与1，圆心距\(|m+2|\)。有公共点\(\Longleftrightarrow|r_{1}-r_{2}|\leqslant|C_{1}C_{2}|\leqslant r_{1}+r_{2}\)，即\(2\leqslant|m+2|\leqslant 4\)，解得\(2\leqslant m+2\leqslant 4\)或\(-4\leqslant m+2\leqslant -2\)，即\(m\in[-6,-4]\cup[0,2]\)。}
\end{ansblock}

\tihao{14}\zu{含参方程表示曲线的分类讨论×解答} \tieside{中档(知识点三)}就实数\(k\)的取值，讨论方程\(x^{2}+2y^{2}-2x-8y+k=0\)表示的曲线形状。
\liubai[16mm]{解答书写区}

\begin{ansblock}[2章-练-课时10-中4]
% ans:2章-练-课时10-中4
\ansitem{14}{k<9：椭圆；k=9：单点(1,2)；k>9：不表示任何曲线}
\ansnote{详解}{配方\((x-1)^{2}+2(y-2)^{2}=9-k\)。\(k<9\)时为椭圆\((x-1)^{2}/(9-k)+(y-2)^{2}/((9-k)/2)=1\)；\(k=9\)时方程表示单点\((1,2)\)（退化椭圆）；\(k>9\)时无实解，不表示任何曲线。}
\end{ansblock}

% ============================================================
% 思维探索（题15~16＝难1~难2：多结论选择＋填空，居末档照库序）
% ============================================================
\dangbadge{思}{维}{探}{索}

\tihao{15}\zu{新定义曲线·性质研究×多结论选择} \tieside{难(知识点三)}双纽线，也称伯努利双纽线，其描述首见于1694年，雅各布·伯努利将其作为椭圆的一种类比来处理．椭圆是由到两个定点距离之和为定值的点的轨迹，而卡西尼卵形线则是由到两定点距离之乘积为定值的点的轨迹，当此定值使得轨迹经过两定点的中点时，轨迹便为伯努利双纽线．伯努利将这种曲线称为lemniscate，为拉丁文中“悬挂的丝带”之意．双纽线在数学曲线领域的地位占有至关重要的地位．双纽线像数字“8”，不仅体现了数学的对称、和谐、简洁、统一的美，同时也具有特殊的有价值的艺术美，是形成其它一些常见的漂亮图案的基石，也是许多设计者设计作品的主要几何元素．曲线\(C\)：\((x^{2}+y^{2})^{2}=4(x^{2}-y^{2})\)是双纽线，则下列结论正确的是\nobreak(\kongwei)〔图注：“8”字形双纽线示意，照录自源〕

\lxopt{A．曲线\(C\)经过5个整点（横、纵坐标均为整数的点）}

\lxopt{B．曲线\(C\)上任意一点到坐标原点\(O\)的距离都不超过2}

\lxopt{C．曲线\(C\)关于直线\(y=x\)对称的曲线方程为\((x^{2}+y^{2})^{2}=4(y^{2}-x^{2})\)}

\lxopt{D．若直线\(y=kx\)与曲线\(C\)只有一个交点，则实数\(k\)的取值范围为\((-\infty,-1]\cup[1,+\infty)\)}

\begin{ansblock}[2章-练-课时10-难1]
% ans:2章-练-课时10-难1
\ansitem{15}{BCD}
\ansnote{详解}{A：令\(y=0\)得\(x^{4}=4x^{2}\)，\(x=0\)或\(\pm 2\)，整点\((2,0)\)、\((-2,0)\)、\((0,0)\)三个；若整点满足\(y\neq 0\)，则由B项结论\(x^{2}+y^{2}\leqslant 4\)只能\(y=\pm 1\)、\(x\in\{0,\pm 1,\pm 2\}\)：代入\((x^{2}+1)^{2}=4(x^{2}-1)\)，记\(X=x^{2}\)得\(X^{2}-2X+5=0\)，判别式\(4-20<0\)无解——\(y\neq 0\)无整点。故曲线仅过3个整点，“5个整点”错。\par\noindent B：由\((x^{2}+y^{2})^{2}=4(x^{2}-y^{2})\leqslant 4(x^{2}+y^{2})\)得\(x^{2}+y^{2}\leqslant 4\)，即\(|OP|\leqslant 2\)。C：方程中\(x\)、\(y\)互换得\((x^{2}+y^{2})^{2}=4(y^{2}-x^{2})\)，即关于\(y=x\)对称的曲线。\par\noindent D：\(y=kx\)代入：\(x^{4}(1+k^{2})^{2}=4x^{2}(1-k^{2})\)。\(x=0\)恒给交点\((0,0)\)；\(|x|>0\)的解须\(x^{2}=4(1-k^{2})/(1+k^{2})^{2}>0\Longleftrightarrow|k|<1\)。故\(|k|\geqslant 1\)时只有原点一个交点；\(|k|<1\)时另有\(\pm\)两交点（\((0,0)\)之外），不止一个（\(k=\pm 1\)时\(x^{4}\cdot 2=0\)仅原点，含入界点正确）。选BCD。}
\end{ansblock}

\tihao{16}\zu{含参曲线系多结论判断×填空} \tieside{难(知识点三)}关于曲线\(C\)：\((x-m)^{2}+(y-2m)^{2}=n^{2}/2\) 有以下五个结论：
①当\(m=1\)时，曲线\(C\)表示圆心为\((1,2)\)，半径为\((\sqrt{2}/2)|n|\)的圆；
②当\(m=0\)，\(n=2\)时，过点\((3,3)\)向曲线\(C\)作切线，切点为\(A\)，\(B\)，则直线\(AB\)的方程为\(3x+3y-2=0\)；
③当\(m=1\)，\(n=\sqrt{2}\)时，过点\((2,0)\)向曲线\(C\)作切线，则切线方程为\(y=-(3/4)(x-2)\)；
④当\(n=m\neq 0\)时，曲线\(C\)表示圆心在直线\(y=2x\)上的圆系，且这些圆的公切线方程为\(y=x\)或\(y=7x\)；
⑤当\(n=4\)，\(m=0\)时，直线\(kx-y+1-2k=0\)（\(k\in R\)）与曲线\(C\)表示的圆相离．
以上正确结论的序号为\kongda{②④}．

\begin{ansblock}[2章-练-课时10-难2]
% ans:2章-练-课时10-难2
\ansitem{16}{②④}
\ansnote{详解}{①\(m=1\)时\((x-1)^{2}+(y-2)^{2}=n^{2}/2\)，但\(n=0\)时退化为点\((1,2)\)，不（恒）表示圆——以偏概全，错误。②\(m=0\)、\(n=2\)：\(C\)：\(x^{2}+y^{2}=2\)。记\(D(3,3)\)，\(|DO|=3\sqrt{2}\)，切线长\(|DA|=|DB|=\sqrt{18-2}=4\)，切点\(A\)、\(B\)即\(C\)与圆\((x-3)^{2}+(y-3)^{2}=16\)（圆心\(D\)半径4）的公共点，公共弦方程两式相减：\(x^{2}+y^{2}-2-(x^{2}+y^{2}-6x-6y+2)=0\)，即\(3x+3y-2=0\)，正确。\par\noindent ③\(m=1\)、\(n=\sqrt{2}\)：\(C\)：\((x-1)^{2}+(y-2)^{2}=1\)。过\((2,0)\)的切线：竖直直线\(x=2\)到圆心\((1,2)\)距离为\(1=\)半径，已是切线；另设\(y=m(x-2)\)：\(|m-2-2m|/\sqrt{m^{2}+1}=1\rightarrow(m+2)^{2}=m^{2}+1\rightarrow m=-3/4\)，即\(3x+4y-6=0\)也是切线。结论只给后者、漏\(x=2\)，“切线方程为”不完备，错误。\par\noindent ④\(n=m\neq 0\)：\((x-m)^{2}+(y-2m)^{2}=m^{2}/2\)，圆心\((m,2m)\)恒在\(y=2x\)上，半径\((\sqrt{2}/2)|m|\)。圆心到\(y=x\)距离\(=|m-2m|/\sqrt{2}=(\sqrt{2}/2)|m|=\)半径；到\(y=7x\)距离\(=|7m-2m|/\sqrt{50}=(\sqrt{2}/2)|m|=\)半径——两直线都是圆系公切线。穷尽性补核（源未给）：设公切线\(y=kx+b\)对一切\(m\)成立，则\([m(k-2)+b]^{2}=\frac{1}{2}m^{2}(k^{2}+1)\)恒成立，比较系数：\(b=0\)且\((k-2)^{2}=(k^{2}+1)/2\)，即\(k^{2}-8k+7=0\)，\(k=1\)或7；竖直线\(x=c\)代入\(|m-c|=(\sqrt{2}/2)|m|\)不能对一切\(m\)成立，无。故公切线恰为\(y=x\)、\(y=7x\)，正确。\par\noindent ⑤\(n=4\)、\(m=0\)：\(C\)：\(x^{2}+y^{2}=8\)，直线\(kx-y+1-2k=0\)即\(k(x-2)-(y-1)=0\)恒过定点\((2,1)\)，而\(4+1=5<8\)，定点在圆内，直线恒与圆相交而非相离，错误。综上，正确的是②④。}
\end{ansblock}

% ---- 尾块（\tailfill 置 \end{multicols} 前；无参＝内容尾随框，每件恰 1 处） ----
\tailfill

\end{multicols}
\end{document}
